\section{Einleitung}

Die Bodenatmung beschreibt die Freisetzung von Kohlenstoffdioxid (CO$_2$) aus dem Boden und bildet den 
zweitgrößten Kohlenstofffluss in terrestrischen Ökosystemen nach der Photosynthese. \cite{raich1992global} Sie setzt
sich aus der mikrobielle Aktivität und der Wurzelatmung zusammen und bildet dafür zentrale Maße. 
Der Anteil der Wurzelatmung schwankt je nach Ökosystem und Jahreszeit zwischen 10\% und 90\%. \cite{Hanson2000}
Die Bodenatmung stellt einen wichtigen Bestandteil des globalen Kohlenstoffkreislaufs dar. Böden sind das
größte terrestrische Reservoir für organischen Kohlenstoff und speichern etwa 71\% bis 81\% des gesamten
terrestrischen Kohlenstoffs. \cite{IPCC2001} Die Bodenatmung ist der Hauptweg auf dem Kohlenstoff wieder in die 
Atmosphäre zurückkehrt, wodurch kleine Änderungen in dessen Intensität große Auswirkungen auf die globale 
CO$_2$-Konzentration und damit auf das Klima haben kann. \cite{1997biogeochemistry,IPCC2001,Churkina2003} \\

Im Nationalpark Hainich wurden umfangreiche Messungen zur Bodenatmung sowie zu physikalischen, chemischen und
biologischen Bodenparametern durchgeführt. Diese Daten ermöglichen die statistische Modellierung der Bodenatmung
in Abhängigkeit von Umweltvariablen wie Bodentemperatur, Bodenfeuchte, pH-Wert und Nährstoffgehalten. Das Ziel
der Forschung ist es, die Sensitivität der Bodenatmung gegenüber dem Klimawandel besser vorherzusagen. \cite{diss} \\

Ziel dieser Arbeit ist die Herleitung eines einfachen, aber aussagekräftigen Prognosemodells für die
Bodenatmung auf Basis des gegebenen Datensatzes. Dies wird mit einem klassischen linearem Modell auf
log-transformierten Daten umgesetzt. Darüber hinaus wird im Rahmen einer Simulationsstudie
untersucht, welche Auswirkungen eine falsche Anwendung von t-Tests im Kontext der Vorwärtsselektion von
Prädiktoren auf die Irrtumswahrscheinlichkeit hat.
